\section{当变量不再是向量：从 Boyd 到函数空间}

\label{sec:0-1}

控制变量是函数而不是向量时，Boyd 为什么不够？更准确地说，Boyd 的有限维母语为什么不能原封不动地搬过来？原因不在于凸性失效了，而在于承载凸性的空间变了。只要未知量从 $\mathbb R^n$ 里的向量变成了一个函数、一个测度，或者一条随机轨道，问题就不再只是“坐标更多”，而是整个极限、连续性和对偶对象都换了舞台。

最短的桥接例子是把有限维二次规划和一个函数空间上的变分问题并排来看。有限维里，我们熟悉
\[
\min_{x\in\mathbb R^n}\left\{\frac12 x^\top Qx+c^\top x\right\}.
\]
如果把未知量换成控制函数 $u\in L^2(0,T)$，就会自然遇到
\[
\min_{u\in L^2(0,T)}
\left\{
\frac12\int_0^T |u(t)|^2\,dt+\Phi\!\left(\int_0^T B(t)u(t)\,dt\right)
\right\}.
\]
这两个问题共享同一层凸优化直觉：目标由一个二次项和一个复合项组成，后者都通过线性映射进入。然而它们的数学成本完全不同。前者的可行集、收敛和乘子都能在矩阵语言中处理；后者则立刻要求我们知道 $L^2(0,T)$ 的拓扑、弱收敛、对偶配对以及算子意义下的最优性条件。

\begin{example}[有限维与无穷维的最短桥]
把上面的函数空间问题投影到一个有限维基底
\[
u(t)\approx \sum_{k=1}^m \alpha_k \varphi_k(t),
\]
就会得到一个关于系数向量 $\alpha\in\mathbb R^m$ 的二次规划。反过来看，Boyd 中那些熟悉的二次规划、线性约束和投影算法，正是更一般函数空间问题在有限维离散化后的坍缩版本。把这个例子放在这里，是为了提醒读者：本书不是离开 Boyd，恰恰是在解释 Boyd 的母语来自哪里。
\end{example}

\paragraph{四个失效点}

一旦变量不再是向量，至少有四件在有限维里默认成立的事不再自动成立。

\begin{enumerate}
  \item \textbf{紧性失效}：闭有界集未必紧，极小化过程可能只得到有界而得不到强收敛极限。
  \item \textbf{极限语言升级}：弱拓扑和弱$^\ast$拓扑往往不是第一可数的，序列不再足以刻画闭包和紧性。
  \item \textbf{导数与乘子升级}：梯度不一定还是向量，最优性条件必须进入对偶空间、次微分和伴随算子。
  \item \textbf{算法接口升级}：矩阵迭代要被重新写成邻近算子、预解式和单调算子分裂。
\end{enumerate}

\begin{remark}
这四个失效点几乎就是本书后文章节安排的理由。拓扑与描述集合论解决“极限怎么说”；测度论与随机空间解决“概率对象怎么放进优化”；泛函分析与对偶理论解决“乘子和支持超平面属于哪里”；单调算子与分裂算法解决“如何把无穷维最优性系统变成可计算更新”。
\end{remark}

因此，本书的“升维”不是炫技，而是对问题本身的诚实。只要你关心的是最优控制、随机策略、图信号、核函数表示或分布式一致性，那么变量就早已不是一个向量；剩下的工作，不过是承认这一点，并为它支付必要的数学代价。第 \ref{sec:0-2} 节接下来要做的，就是把这笔代价拆开说明：为什么必须上拓扑、弱收敛、对偶与算子。Boyd 在这里仍然是重要参照，但他更像一条有限维影子，而不是终点。
